\section{Risk Analysis}

\label{sec:risks}

\subsection{Taxonomy of Failure Modes}
\label{sec:failmodes}

We enumerate the principal failure modes and bound the damage each can cause
under HHP constraints.

\paragraph{Risk R1: A generated tool makes things worse.}
A Tier 1 tool that regresses performance (e.g., incorrectly implements a
function, introduces latency) could harm task completion during its TTL window.

\textbf{Bound:} The damage is bounded to the TTL window (7 days) and to tasks
matching the tool's trigger pattern. By the validation gate $\Gamma_1$, the
tool must pass all known triggering failure cases before deployment. Regression
on \emph{previously passing} cases is impossible if the trigger set is
comprehensive. Regression on \emph{novel} cases is bounded by the tool's
invocation rate over 7 days. Telemetry monitoring (Section~\ref{sec:cycle})
detects regressions within $O(\text{invocations})$ time and the tool can be
manually deregistered before TTL expiry.

\paragraph{Risk R2: An agent modifies security-sensitive code.}
A Tier 4 modification targeting authentication, key management, or cryptographic
code could introduce vulnerabilities.

\textbf{Bound:} By the Worktree Isolation Theorem (Theorem~\ref{thm:isolation}),
the modification cannot affect the running system until merged. The Tier 4 gate
requires both CI (which typically includes security scans) and human review.
Human reviewers are expected to apply heightened scrutiny to security-sensitive
paths. In practice, repositories containing security-sensitive code (e.g., the
KMS system) should be added to an explicit exclusion list $\mathcal{E}_{\text{excl}}$
such that $\text{ClassifyScope}(\mathcal{F})$ returns $k = \text{forbidden}$ for
any friction event involving $\mathcal{R} \in \mathcal{E}_{\text{excl}}$.

\paragraph{Risk R3: Two agents modify the same component concurrently.}
Two agents independently identifying friction in the same component and
submitting conflicting Tier 4 PRs.

\textbf{Bound:} Git merge conflicts are detected at PR merge time, not at
worktree creation time. The conflict is surfaced to the human reviewer who
can choose which change to apply or request a combined modification. Since
worktrees are isolated (Theorem~\ref{thm:isolation}), conflicting worktrees
cannot corrupt each other; only the merge resolution step requires intervention.
In multi-agent deployments, a distributed lock on the target repository path
can serialize Tier 4 hack cycles to prevent concurrent conflicts.

\paragraph{Risk R4: The modification loop never terminates.}
An agent in the MODIFY--TEST loop that cannot produce a passing modification
could loop indefinitely.

\textbf{Bound:} By Invariants 2 and 3, each hack cycle is bounded at 3 fix
attempts and $T_{\text{hack}} = 5$ minutes. After 3 failed attempts, the loop
terminates, the friction event is archived, and the session budget is
\emph{not} decremented (an abandoned hack does not consume the session budget).
The time-box provides an absolute wall-clock bound independent of fix attempt
count.

\paragraph{Risk R5: An agent promotes a harmful tool to a higher tier.}
A Tier 1 tool that produces subtly wrong outputs but meets the basic validation
criteria ($\Gamma_1$) could be promoted to Tier 2 with a broader invocation scope.

\textbf{Bound:} Promotion from Tier 1 to Tier 2 requires $\tau(\rho, t) > 0.8$
after $\geq 10$ uses. This means the tool must have produced successful outcomes
in at least 80\% of its invocations over a statistically significant number of
uses. A subtly wrong tool will manifest in the telemetry (Section~\ref{sec:cycle})
as reduced task completion rates, which will suppress its trust score below the
promotion threshold.

\subsection{Comparative Risk Table}
\label{sec:risk-table}

\begin{table}[H]
\centering
\begin{tabularx}{\textwidth}{l X X c}
\toprule
Risk & Failure Mode & HHP Mitigation & Max Blast Radius \\
\midrule
R1 & Tool regresses performance & TTL + telemetry monitoring & 7 days, tool-scoped \\
R2 & Security-sensitive modification & Exclusion list + human review & Zero (gated) \\
R3 & Concurrent conflicting PRs & Git conflict detection + reviewer & PR-scoped \\
R4 & Non-terminating loop & 3-attempt limit + 5-min time-box & 5 min session overhead \\
R5 & Harmful tool promotion & Trust score threshold (0.8, 10 uses) & Tier-bounded \\
\bottomrule
\end{tabularx}
\caption{Risk analysis with HHP mitigations and maximum blast radius under
the protocol constraints.}
\label{tab:risks}
\end{table}

\subsection{Adversarial Threat Model}
\label{sec:adversarial}

We briefly consider an adversarial setting where an attacker can influence the
agent's friction event classification (e.g., via prompt injection in tool outputs).

\begin{proposition}[Adversarial Containment]
\label{prop:adversarial}
Under the HHP constraints, an adversary who can inject arbitrary friction events
into the agent's context cannot deploy a modification that bypasses at least
one human gate (for Tier 4) or causes system-wide impact (for Tiers 1--3).
\end{proposition}

\begin{proof}
For Tier 4: any modification requires $\Gamma_4 = $ CI $\wedge$ HumanApproval.
An adversary cannot forge a human approval without compromising the human
reviewer or the PR platform. For Tiers 1--3: modifications are scoped to the
individual agent's session or extension registry; they cannot affect shared
infrastructure. Tier 1 tools expire after 7 days. Thus, the maximum persistence
of any adversarially-injected modification is bounded. \qed
\end{proof}

% ============================================================================
% 9. DISCUSSION AND FUTURE WORK
% ============================================================================
